\SourceLeaf{420}%
offenbar
\[
\begin{aligned}
\delta x&=\delta\varrho.\cos(\varphi+90^{0})\\
\delta y&=\delta\varrho.\sin(\varphi+90^{0})
\end{aligned}
\]
oder
\[
\begin{aligned}
\delta x&=-\delta\varrho.\sin\varphi\\
\delta y&=+\delta\varrho.\cos\varphi.
\end{aligned}
\]
Wenn $\delta\varrho$ unendlich klein ist, gilt
\[
\begin{aligned}
\delta V&=p\delta x+q\delta y\\
&=(-p\sin\varphi+q\cos\varphi)\delta\varrho\\
&=\mp\delta\varrho\BareRoot{(pp+qq)}
\end{aligned}
\]
und auf der Kurve selbst ist $V$ $=0$. Daher gelten die oberen Vorzeichen, wenn $V$ beim Durchgang durch die Kurve von links nach rechts vom Positiven ins Negative übergeht, und umgekehrt. Verbindet man dies mit dem, was am Ende des 2. Artikels gesagt wurde, so folgt, dass die Kurve immer nach der Seite konvex ist, auf der $V$ das Vorzeichen annimmt, das
\[
Pqq-2Qpq+Rpp
\]
hat.

Ist z. B. die Kurve ein Kreis und setzt man
\[
V=xx+yy-aa,
\]
so gilt
\[
\begin{gathered}
p=2x,\qquad q=2y,\\
P=2,\qquad Q=0,\qquad R=2,\\
Pqq-2Qpq+Rpp=8yy+8xx=8aa,\\
(pp+qq)^{\frac32}=8a^3,\\
r=\pm a
\end{gathered}
\]
und die Kurve ist nach der Seite konvex, auf der
\[
xx+yy>aa,
\]
wie es sein soll.
